\documentclass[11pt,a4paper]{article}

% ------------------------------------------------------------------ setup
\usepackage[margin=2.4cm,top=2.2cm,bottom=2.6cm]{geometry}
\usepackage{graphicx}
\usepackage{booktabs}
\usepackage{float}
\usepackage{xcolor}
\usepackage{parskip}
\usepackage{microtype}
\usepackage{caption}
\usepackage{fancyhdr}
% Keep every figure and table inside the section that introduces it: a
% barrier before each section flushes any pending floats first, so diagrams
% never drift under the wrong heading.
\usepackage[section]{placeins}
\usepackage{hyperref}

\definecolor{accent}{RGB}{2,132,199}
\definecolor{dimtext}{RGB}{100,116,139}
\definecolor{darktext}{RGB}{30,41,59}

% Internal references in plain dark text, not blue; only real URLs stay colored.
\hypersetup{colorlinks=true, linkcolor=darktext, urlcolor=accent}
\captionsetup{font=small, labelfont={bf,color=darktext}, skip=8pt}

% Let figures float gracefully instead of forcing them or hoarding them.
\renewcommand{\topfraction}{0.85}
\renewcommand{\bottomfraction}{0.75}
\renewcommand{\textfraction}{0.15}
\renewcommand{\floatpagefraction}{0.75}
\setcounter{topnumber}{3}
\setcounter{bottomnumber}{2}
\setcounter{totalnumber}{4}

\makeatletter
\setlength{\@fptop}{0pt}
\setlength{\@fpsep}{20pt}
\setlength{\@fpbot}{0pt plus 1fil}
\makeatother

\pagestyle{fancy}
\fancyhf{}
\fancyhead[L]{\small\color{dimtext} Hyrox: Every Race Ever Run}
\fancyhead[R]{\small\color{dimtext} \thepage}
\renewcommand{\headrulewidth}{0.3pt}

\graphicspath{{../analysis/output/figures/}}

% Figure: full page width, capped height so it fits between text.
% The first argument can be used to override the default key-value options for \includegraphics.
\newcommand{\chart}[3][width=\linewidth,height=0.42\textheight]{%
\begin{figure}[!htbp]\centering
\includegraphics[#1,keepaspectratio]{#2.png}
\caption{#3}\label{fig:#2}
\end{figure}}

% Table fragment from tables/
\newcommand{\tabfrag}[2]{%
\begin{table}[!htbp]\centering\small
\input{tables/#1}
\caption{#2}\label{tab:#1}
\end{table}}

\title{\vspace{-1em}{\Huge\bfseries Hyrox: Every Race Ever Run}\\[0.6em]
{\Large What 1.18 million finish lines tell us about the world's fastest-growing fitness race}}
\author{{\large\bfseries Devansh Khetan}\\[0.3em]
{\color{dimtext} An analysis of the complete Hyrox results archive, Seasons 1 to 9 (2018 to 2026)}}
\date{\color{dimtext} July 2026}

\begin{document}

\maketitle
\thispagestyle{empty}

\begin{abstract}
\noindent
In the winter of 2017, a few thousand curious athletes turned up at exhibition
halls in nine German-speaking cities to try a new kind of fitness race. Eight
seasons later that race fills arenas on five continents, and it has quietly
produced one of the richest datasets in amateur sport: every athlete, every
kilometre run, every SkiErg, sled and wall ball, timed and published. I worked
through all of it, 1{,}181{,}411 deduplicated results from 145 host cities,
to answer the questions every Hyrox athlete eventually asks. How fast is fast?
Where does the race actually hurt? Do negative splits work? What separates
podium athletes from the rest of us? The short version: Hyrox is a running
race wearing a strongman costume, pacing discipline is worth minutes, balance
beats brilliance, and almost every all-time great performance was set in the
last twenty-four months. A note before we begin: everything here is built on
results data that was publicly scrapable from Hyrox's servers, so it inherits
whatever gaps and quirks that archive contains.
\end{abstract}

\vspace{1em}
\tableofcontents
\newpage
\listoffigures
\listoftables
\newpage

% ================================================================== 1
\section{Introduction: the race, the first season, and this data}

If you have never stood in a Hyrox arena, the format takes one sentence to
explain and about ninety minutes to survive. You run one kilometre, then do a
workout station, and you repeat that eight times: a thousand metres on the
SkiErg, a 50 metre sled push, a 50 metre sled pull, 80 metres of burpee broad
jumps, a thousand metres on the rowing machine, a 200 metre farmers carry,
100 metres of walking lunges with a sandbag, and finally one hundred wall
balls. Between each run and station you pass through a transition area the
organisers call the roxzone, and the clock never stops. Every venue in the
world uses the same distances, the same weights and the same layout, which is
the quiet genius of the whole thing: a Hyrox time in Mumbai means the same as
a Hyrox time in Manchester, so everyone on earth is effectively in one long
race.

The first season was a modest, almost entirely German affair. Nine cities
hosted races in the winter of 2017/18: Hamburg was the biggest with about
1{,}500 entries, followed by Karlsruhe, Hannover, N\"urnberg, Stuttgart,
Essen, Oberhausen and Leipzig, with Vienna the only stop outside Germany. In
total, 5{,}254 results were recorded that first winter. The fastest man of the
season was Lukas Storath, who ran 58:39 in Stuttgart; the fastest woman was
Nele Laing with 1:07:51 in Karlsruhe. Both marks would still be respectable
today, which says something about how well-designed the race was from day
one. Eight seasons later a single season produces close to six hundred
thousand results, the men's all-time best has fallen to 50:38 (Alexander
Roncevic, Cologne) and the women's to 56:53 (Lena Putters, Maastricht), and
athletes of more than two hundred nationalities have crossed a Hyrox finish
line.

That growth left behind a remarkable paper trail. Hyrox publishes complete
results with full split times: all eight runs, all eight stations, the
roxzone, ranks and demographics. I collected the entire archive, about 1.67
million raw rows, and cleaned it before drawing a single chart. The cleaning
mattered more than usual. Multi-day events publish the same athletes twice,
once in a day listing and again in an overall listing, and removing those
duplicates cut the archive by nearly a third, down to 1{,}181{,}411 unique
results. The 127 messy division labels in the raw data collapse into ten real
formats, of which four matter most in this report: Open (the standard solo
race), Pro (heavier weights), Doubles (two athletes share the station work)
and Pro Doubles. And because the archive has no stable athlete identifier, I
track people across races by name and sex, which works well at scale but
makes career numbers estimates rather than exact counts. By that measure
roughly 906{,}000 distinct athletes appear in the data. For anything about
performance I use only what I call clean races: a plausible finish time with
all sixteen splits present, which about 929{,}000 results satisfy. Many of
the doubles entries carry only a finishing time and no station splits, so
they count toward participation totals but never toward the performance
analysis, which stays anchored on fully-timed solo races.

One honest warning before the numbers start. This entire report is based on
data that was scrapable from Hyrox's public results servers. It is the most
complete picture of the sport available outside the company, but it is not an
official export: athletes who never finished are largely missing, some fields
are patchy in early seasons, and nothing here has been verified by Hyrox
itself. Where a quirk of the data affects a conclusion, I say so in the text.

% ================================================================== 2
\section{From nine German cities to the world}

The simplest way to see what happened to Hyrox is Table~\ref{tab:season_summary}.
Every season roughly doubled or better after the pandemic, and the sport that
managed 5{,}254 results in its first winter closed Season 8 with 679{,}500.
One clarification, because the table can look odd at first glance: the
nationalities column counts where the \emph{athletes} come from, not where
races are held, which is why it exceeds the host city count. Races have been
hosted in around thirty countries so far, but people from 235 nations have
travelled to start one. Season 9 had barely begun when I collected the data,
so its row (and its points in several charts later) should be read as a few
opening weekends, not a collapse.

\tabfrag{season_summary}{Season by season vital signs. Athlete counts are
name-based estimates, and the women's share is measured among individually
sexed athletes (mixed-gender doubles teams are recorded separately and
excluded from that column). Season 3 collapsed to a handful of COVID-era
events, and Season 9 had just begun when the data was collected.}

Figure~\ref{fig:01_participants_per_year} shows the same story by calendar
year, and the shape is worth pausing on. The 2020 hole is COVID, obviously,
but the interesting part is what happens afterwards: the recovery does not
rejoin the old trend line, it launches off a much steeper one. Growth driven
by a fad usually resets after an interruption. Growth that accelerates
through an interruption is usually structural. Note also how closely the
unique-athletes line tracks the results bars, which tells us the boom comes
from new people arriving rather than the same people racing more often. The
2026 bar, like Season 9 in the table, is still filling in.

\chart{01_participants_per_year}{Race results per calendar year, with the
number of distinct athletes behind them. The 2026 bar covers only the opening
weeks of Season 9.}

The supply side of this growth is the race calendar itself, and it expanded
just as aggressively. Figure~\ref{fig:01_cities_per_season} counts host
cities per season: nine in the first year, twelve in the second, then after
the pandemic 23, 35, 50, 66 and finally 102 in Season 8. Every one of those
cities needs a hall, tonnes of identical equipment and a local community
willing to show up, so this curve is the hardest part of the business to
fake. Figure~\ref{fig:01_top_cities} shows where the accumulated volume
lives: London leads the all-time list comfortably, thanks to multiple events
per season and the giant Olympia race weeks, with Germany's original venues
and the big American cities filling out the list.

\chart{01_cities_per_season}{Distinct host cities per season.}

\chart{01_top_cities}{The 25 biggest host cities by all-time results. London
has hosted more Hyrox results than any city on earth.}

Figure~\ref{fig:01_city_season_heatmap} adds the time dimension, and it is my
favourite way to see how young this sport really is. Each row is one of the
top twenty cities and each column a season, with brightness showing volume.
Almost every row is dark until Season 5 or 6 and then suddenly lights up.
Only Hamburg and a couple of other German venues have any meaningful history
before the pandemic. Most of what we now think of as the sport's flagship
cities have hosted Hyrox for barely three years.

\chart[width=0.9\linewidth,height=0.58\textheight]{01_city_season_heatmap}{When each big city joined the calendar. Most
of today's flagship venues only appear from Season 5 onwards.}

Put it all on a map, Figure~\ref{fig:01_world_map}, and the honest summary is
that Hyrox is a Europe-and-America sport that has started reaching further.
The colour is deepest across Germany, the UK and the US, with Australia
strong for its size. The Middle East has a visible foothold, Africa has a
first one through South Africa, and Asia's presence is young but, as we will
see shortly, growing at a pace nowhere else can match. Hyrox is global in
ambition and increasingly global in footprint, but the volume still lives
overwhelmingly in Europe and North America. The blank patches on this map are
the next decade of the sport.

\chart[width=\linewidth,height=0.55\textheight]{01_world_map}{All-time results by host country. Europe and the US
dominate today; the Middle East, Africa and Asia are the expansion frontier.}

That frontier is already visible in the data if you look at athletes rather
than venues. Figure~\ref{fig:01_emerging_markets} tracks results by athlete
nationality for five developing markets on a logarithmic scale, which is the
only scale that can hold them. India went from about 100 results in Season 6
to 794 in Season 7 to 5{,}861 in Season 8. China went from 65 to more than
8{,}000 over the same window; South Africa from 123 to about 5{,}500; Mexico,
the most established of the group, passed 15{,}000. These are growth rates of
roughly ten times per season, from admittedly small bases. Germany took eight
seasons to build its volume. If even half of these slopes hold, the sport's
centre of gravity will look very different in three years, and countries like
India and China are plausible future superpowers of the sport.

\chart{01_emerging_markets}{Results by athlete nationality for five emerging
markets, on a log scale. India, China and South Africa are compounding at
roughly ten times per season.}

% ================================================================== 3
\section{The formats, and how Doubles took over}

Hyrox is not one race but a family of them, and the balance between the
formats has shifted in a way I think is the single most important trend in
the participation data. Figure~\ref{fig:01_participation_by_season_division}
stacks each season's results by division. For years Open, the classic solo
race, was the biggest block, but watch the green band overtake it: Doubles,
where two athletes run every lap together and split the station work between
them, has grown faster than the solo formats in every post-COVID season, and
in Season 8 it passed Open outright. Doubles alone drew 305{,}731 Season 8
results against Open's 268{,}055, and together with Pro Doubles the
partnered formats now make up 51 percent of the sport. The majority format
in Hyrox is no longer the individual race. Hyrox has quietly become a team
sport.

\chart{01_participation_by_season_division}{Season results stacked by
division. Doubles overtook Open in Season 8, and the partnered formats now
make up 51 percent of all results.}

Figure~\ref{fig:01_division_growth_rate} makes the comparison explicit. The
left panel shows each division's absolute size; the right panel shows the
season-on-season growth rate, which strips out Open's head start. I start
the growth panel at Season 5 deliberately: the COVID collapse and rebound
produce meaningless four-digit percentages in Seasons 3 and 4 that would
squash everything else in the chart into a flat line at the bottom. From
Season 5 onwards, on a clean comparison, the doubles curves sit above the
solo curves essentially throughout.

\chart{01_division_growth_rate}{Division size (left) and season-on-season
growth from Season 5 onwards (right). The COVID rebound seasons are excluded
from the growth panel because their percentages are meaninglessly large.}

Why does Doubles recruit so well? The clearest answer in the data comes from
comparing what a partner actually changes, shown in
Figure~\ref{fig:08_doubles_vs_solo}. The running is untouched: doubles pairs
run every kilometre together, and their median run total (42:06 for men) is
almost identical to solo runners' (42:15). The stations are transformed:
men's wall balls drop from a median of 6:50 solo to 4:18 in doubles, and
every station shows the same halving pattern, because partners split the
reps. A doubles race is a solo race with half the strength work, which
removes exactly the part that intimidates newcomers, and you get to bring a
friend. It is not hard to see why it is eating the sport, and its social
nature may also be part of why the field keeps getting more balanced between
men and women, which is where we turn next.

\chart{08_doubles_vs_solo}{Median segment times, solo versus doubles. Running
is identical; every station roughly halves. The entire doubles advantage
comes from shared station work.}

% ================================================================== 4
\section{Who races: gender, age and nations}

Hyrox likes to describe itself as a sport for everybody, and the demographics
increasingly back that up. Start with gender,
Figure~\ref{fig:01_sex_split}. Men still outnumber women, but the share of
women has risen every single season without exception: 33 percent in Season
1, 42 percent by Season 8, and the early Season 9 data sits at 45 percent.
Very few strength-flavoured sports get anywhere near parity, and Hyrox is
visibly on its way. I suspect the doubles boom is part of the mechanism,
since mixed and women's pairs lower the entry barrier in exactly the way the
previous section described, and the standardized, scalable format means women
compete on the same course with appropriately scaled implements rather than
in a separate event. Whatever the exact mix of causes, the trend line is one
of the healthiest in the whole dataset.

\chart{01_sex_split}{Results by sex (left) and women's share of the field
(right). The share has risen every season, from 33 to 45 percent.}

Age tells a similar story of breadth, with one nuance worth dwelling on.
Figure~\ref{fig:01_age_distribution} shows the all-time age structure: the
sport peaks at 30 to 34, the 25 to 39 range dominates the field, and yet the
masters brackets are enormous, with athletes aged 40 and up contributing
more than a quarter of all results. Part of the young tilt is simply where
Hyrox has recently arrived. The newest markets, India among them, skew younger than
the mature German and British fields, so as those markets grow you would
expect the sport's age mix to tilt further toward the twenties.

\chart{01_age_distribution}{All-time results by age group and sex. The 25 to
39 brackets dominate, but more than a quarter of the field is 40 or older.}

To check whether that tilt is actually happening yet, I tracked the age mix
season by season, Figure~\ref{fig:01_age_growth_by_season}. The answer so
far is: only slightly. The 25 to 34 core holds around the same share in
every season, the 16 to 24 slice edges upward a little, and the 45 plus
share drifts gently down. The mix is remarkably stable for a sport whose
size has grown thirtyfold over the same window, which tells you the new
markets are reproducing the sport's demographic shape rather than changing
it. If the emerging-market growth from the previous section continues, this
is the chart where I would expect the youth shift to show up first.

\chart{01_age_growth_by_season}{Share of each season's results by age band.
The mix is stable so far, with a slight tilt toward the youngest band as new
markets join.}

Nationality completes the picture. Figure~\ref{fig:01_top_nationalities}
ranks the twenty largest Hyrox nations by athlete count, and the top of the
list is exactly who you would expect: Germany, the United Kingdom and the
United States, followed by the Netherlands, France and Australia.
Figure~\ref{fig:01_nationality_growth} adds the time axis and shows the
handover: Germany built the sport almost alone through the first two
seasons, then the British and American curves went vertical after COVID and
overtook. The sport was invented in Germany, but its growth engine now
speaks English.

\chart{01_top_nationalities}{The twenty biggest Hyrox nations by all-time
athlete results.}

\chart{01_nationality_growth}{Results per season for the eight biggest
nations. Germany started it; the UK and US now drive the volume.}

% ================================================================== 5
\section{What a finish looks like}

Ask a marathoner what a good time is and they will give you a number.
Hyrox is young enough that most of its athletes genuinely do not know, and
social media, which mostly amplifies the fastest one percent, actively
miscalibrates them. So let me put the real distribution on the table.
Figure~\ref{fig:02_finish_distribution_open} shows every Open singles finish
in history. The men's median is 1:27:20 and the women's 1:34:46, a gap of
about seven and a half minutes. Half of all men finish between 1:17:38 and
1:39:57, half of all women between 1:24:54 and 1:48:05. If you finish a
Hyrox anywhere around ninety minutes you are, statistically, a perfectly
normal athlete in this sport. Notice also the long right tail with no bump
of stragglers at the end: people having a bad day slow down and grind it
out rather than blowing up, because the format always lets you walk the
next run and breathe at the next station.

\chart{02_finish_distribution_open}{All Open singles finishes ever. Men's
median 1:27:20, women's 1:34:46, with long right tails and no cliff.}

Table~\ref{tab:finish_percentiles} turns that distribution into a lookup:
find your division, your sex and your time, and read off your all-time
percentile. Two things in it deserve comment. Doubles is the fastest
division for the reasons of Section 3. And Pro, the division with the
heavier weights, is \emph{slower} than Open despite attracting stronger
athletes, because the extra sled and carry weight costs more time than the
fitter field gives back. A Pro time and an Open time are simply different
currencies.

\tabfrag{finish_percentiles}{Finish time percentiles by division and sex,
all time. Find your time, read your percentile.}

The percentile curve in Figure~\ref{fig:02_percentile_curve} is the
continuous version, and its shape rewards a moment of study. Through the
broad middle of the field the curve is almost flat, meaning huge crowds of
athletes are separated by a minute or two, so a small improvement moves you
past a lot of people. At the front the curve dives steeply, meaning the
final percentiles cost enormous amounts of fitness to climb. Improving from
the 50th to the 25th percentile takes a man from roughly 1:27 to 1:18, which
is a season of honest training rather than a fantasy.

\chart{02_percentile_curve}{Finish time at every percentile of the all-time
Open field. Flat through the middle, brutal at the front.}

The community's folk benchmarks look different against this backdrop.
Figures~\ref{fig:02_benchmark_histogram_m} and
\ref{fig:02_benchmark_histogram_w} mark them on the men's and women's
distributions, each label showing the share of the all-time field that beats
it. Sub-60 for men, the number the sport's culture treats as the badge of
seriousness, has been achieved in half of one percent of all men's Open
finishes ever. Sub-70 for women is comparably exclusive at 1.7 percent. Even
sub-80 for men already puts you in the top third. These numbers are worth
internalising: the benchmarks are genuinely elite, and a personal best is a
far more sensible target for almost everyone than a number borrowed from the
fastest sliver of the sport.

\chart{02_benchmark_histogram_m}{Men's finishes against the classic
benchmarks. Sub-60 is a top 0.5 percent performance all-time.}

\chart{02_benchmark_histogram_w}{Women's finishes against their benchmarks.
Sub-70 is a top 1.7 percent performance.}

\subsection{Is the field getting faster or slower?}

Here we reach the most commonly misread statistic in Hyrox. If you plot the
median finish per season, Figure~\ref{fig:02_median_finish_trend}, it drifts
slower in every division, and the naive conclusion is that the sport's
athletes are getting worse. They are not. The median is being dragged by
composition: every season, hundreds of thousands of first-timers join, and
first-timers finish mid-pack and behind. When newcomers outnumber veterans
this heavily, the median measures the growth of the base, not the fitness of
any individual. The median is up simply because far more people are
participating.

\chart{02_median_finish_trend}{Median finish per season, by division and
sex. The drift upward is a composition effect of explosive growth, not
athletes getting worse.}

To see what is really happening you have to compare like with like.
Figure~\ref{fig:02_percentile_by_season} overlays the full men's percentile
curve for several seasons, and instead of shifting up or down, the curve
\emph{twists}. At the fast end, recent seasons sit below older ones: the
front 10 percent of the field is faster than it has ever been. At the slow
end recent seasons sit above: the back half is bigger and slower.
Figure~\ref{fig:02_field_depth_trend} shows the same thing as three simple
lines, the 25th, 50th and 75th percentile per season, fanning steadily
apart. The field is not slowing, it is stretching, exactly the way the
marathon stretched when it went from an elite pursuit to a bucket-list
ambition. Both halves of that sentence are signs of health.

\chart{02_percentile_by_season}{Men's Open percentile curves by season. The
curve twists: faster at the front, longer and slower at the back.}

\chart{02_field_depth_trend}{The 25th, 50th and 75th percentile finish per
season. The field is stretching, not slowing.}

\subsection{Age and geography}

Two last cuts of the finish data, because both overturn assumptions. First,
age. Figure~\ref{fig:02_age_curves} plots the median finish by age group,
and the penalty of ageing is far gentler than most people guess: a median
man of 50 to 54 finishes in 1:33:32, only about eight and a half minutes
behind the median 25 to 29 year old's 1:25:03. The curve barely rises
through the thirties and forties. Hyrox rewards aerobic endurance and
discipline, which age slowly and respond to training at any age, far more
than raw speed and power, which do not. For the sport's huge masters field
this is as encouraging as data gets: deep into your fifties, your training
matters more than your birth year.

\chart{02_age_curves}{Median finish by age group with the middle 50 percent
band. A 50 to 54 year old median man concedes only about 8.5 minutes to a
25 to 29 year old.}

Second, geography. Figure~\ref{fig:02_city_difficulty} ranks the fastest and
slowest big venues by median finish, and the spread is enormous, roughly 45
minutes between London Olympia's championship-hardened 1:15 median and
Bengaluru's 2:00. Course layout, heat and altitude all contribute, but the
dominant force is simply who shows up: a qualifier venue attracts trained
veterans while a brand-new market is almost entirely first-timers. This
chart describes fields more than courses, a distinction that matters again
when we visit India at the end of the report.

\chart{02_city_difficulty}{Fastest and slowest big venues by median Open
finish. The 45 minute spread mostly reflects who races there, not course
difficulty.}

% ================================================================== 6
\section{The eight stations}

Between the runs sit the eight stations, and the first thing the data does
is puncture some reputations. Figure~\ref{fig:03_station_medians} shows the
median time each one costs. Wall balls are the biggest bill at 6:50 for the
median man, with the sandbag lunges (5:14) and burpee broad jumps (5:31)
close behind. The sled push, the station everyone fears in advance, takes
just three minutes, and the farmers carry barely two. Fear and cost are not
the same thing: the stations that actually eat your race are the long
grinds in its second half, not the dramatic heavy pushes in its first.

\chart{03_station_medians}{Median time per station in race order. Wall balls
and lunges cost the most; the feared sled push is one of the shortest.}

Zoom out and Figure~\ref{fig:03_time_budget} shows where a whole median race
goes: 49 percent running, 43 percent across all eight stations combined,
and 8 percent in the roxzone, the transition zone between them. Read that
again, because it is the sentence around which this whole report turns.
Running is a bigger share of a Hyrox than every station put together, and
the transitions alone rival any single station. Anyone who trains Hyrox
primarily as a strength event is optimising the smaller half of their race.

\chart{03_time_budget}{Where the median race goes: about half running, 43
percent stations, 8 percent transitions.}

The stations also differ wildly in how much they separate athletes.
Figure~\ref{fig:03_station_spread} shows each station's full spread, and
Figure~\ref{fig:03_station_variability} condenses it to a single relative
number. The erg machines are the great levellers: on the row and SkiErg the
field varies by only 11 to 12 percent, because a rowing machine has a
mechanical ceiling that even strong athletes cannot push far past. Wall
balls sit at the other extreme with 41 percent variation, and the slowest
twentieth of the field needs more than triple the time of the fastest
twentieth there. The practical reading: you can neither win nor lose much
on an erg, so pace them steadily and save your heroics; the sleds, lunges
and wall balls are where races actually move.

\chart{03_station_spread}{Each station's spread, from the 5th to the 95th
percentile. Wall balls stretch the field more than anything else.}

\chart{03_station_variability}{Relative variability by station. Ergs level
the field; wall balls and sleds divide it.}

Who loses time where? Figure~\ref{fig:03_elite_vs_field_gap} compares the
middle and the back of the field against the top five percent, segment by
segment, and the answer is emphatic: the gaps are smallest on the runs and
largest on the sleds, the wall balls and the roxzone. Nearly everyone can
shuffle a kilometre at a survivable pace. What breaks under fatigue and
inexperience is strength-endurance and transition discipline. If you are
mid-pack, this chart is your training plan in one picture.

\chart{03_elite_vs_field_gap}{How much slower the middle and back of the
field are than the top five percent on each segment. The runs separate
least; sleds, wall balls and transitions separate most.}

The comparison between men and women, Figure~\ref{fig:03_sex_gap_by_station},
needs careful reading because Hyrox deliberately scales several implements
by sex, and I have drawn the chart so that the direction is unmissable.
Green bars pointing left mean women are outright faster: wall balls (women
throw a lighter ball to a lower target), the lunges and the sled push, all
scaled stations. Pink bars pointing right mean men are faster, and the
biggest of those, the sled pull and burpee broad jumps, reflect the
strength and power differences the scaling does not fully offset. The
fairest like-for-like comparison is the identical equipment: on the ergs
the gap is a modest 13 to 16 percent, and on the runs about 13 percent. The
chart is really a map of where the sport's scaling equalises the sexes by
design, and it does so remarkably well.

\chart{03_sex_gap_by_station}{Women's median versus men's per segment. Green
means women are faster (scaled implements); the zero line is parity.}

Finally, Figure~\ref{fig:03_station_trend_by_season} tracks each station's
median over the seasons, and every one drifts slower. By now you can guess
why without being told: it is the same first-timer flood that moved the
overall median, arriving at every station simultaneously. If your own
splits are improving against this tide, you are making genuine progress.

\chart{03_station_trend_by_season}{Median men's station times by season. The
drift is the newcomer effect again, visible at every station at once.}

% ================================================================== 7
\section{Pacing: where races are actually decided}

Strip away the equipment and Hyrox asks one brutally simple question: can
you run eight kilometres at an even pace while your heart rate lives
somewhere it should not? Most people cannot, and the data on this is the
most consistent in the entire archive.

Figure~\ref{fig:04_fatigue_curves} is the centrepiece. It plots the median
lap time across the eight runs for four groups of finishers, from the top
five percent of each event to the back twenty. Everyone slows down. The
difference is how much: the top five percent run their final lap 18 percent
slower than their first, while the back of the field runs it 48 percent
slower, finishing at barely half their opening speed. Two athletes with
identical fresh kilometre times can finish twenty minutes apart purely on
how gracefully they decay. Durability, not raw speed, is the quantity this
sport actually tests, and Table~\ref{tab:run_laps_cohort} shows the raw lap
times behind the curves. Look for the kink at lap five, straight after the
burpee broad jumps, and the jump at lap eight, straight after the lunges:
the two stations that most destroy running legs leave fingerprints in
every cohort's laps, elite and novice alike.

\chart{04_fatigue_curves}{Median lap times by finishing group. The top five
percent fade 18 percent from first lap to last; the back of the field fades
48 percent.}

\tabfrag{run_laps_cohort}{The lap times behind the fatigue curves. Note the
kink after the burpee broad jumps (lap 5) and the jump after the lunges
(lap 8).}

Where do athletes run their fastest and slowest kilometres?
Figures~\ref{fig:04_best_lap_position} and \ref{fig:04_worst_lap_position}
answer bluntly. For 57 percent of the field, lap one is the fastest lap of
the entire day, and for 64 percent, lap eight is the slowest. The first
number is a warning dressed as a statistic: a fastest opening lap means you
arrived at the SkiErg having already spent more than you would ever get
back. Fresh legs and start-line adrenaline make that lap feel free. It is
the single most expensive kilometre in the sport.

\chart{04_best_lap_position}{Which lap was each athlete's fastest. For 57
percent of the field, it is lap one, which is a pacing warning, not an
achievement.}

\chart{04_worst_lap_position}{Which lap was each athlete's slowest. For 64
percent it is lap eight, straight off the sandbag lunges.}

The cost of that opening enthusiasm is measurable.
Figure~\ref{fig:04_degradation_vs_percentile} plots each athlete's
last-lap fade against where they finished in their event, and the line is
almost perfectly straight: more fade, worse finish, all the way down the
field. Figure~\ref{fig:04_consistency_vs_percentile} shows the same
discipline from another angle, lap-to-lap steadiness. The best finishers
run like metronomes, varying five to six percent across all eight laps,
while the back of the field is two to three times as erratic. Surging and
recovering is metabolically expensive; even effort is a free efficiency
gain available to anyone willing to run to a plan instead of a feeling.

\chart{04_degradation_vs_percentile}{Last-lap fade against finishing
percentile: a nearly straight line. Going out too hard is the sport's most
measurable mistake.}

\chart{04_consistency_vs_percentile}{Lap-to-lap consistency against results.
Winners vary their laps by five to six percent; the back of the field is far
more erratic.}

The strongest single piece of evidence is the negative split,
Figure~\ref{fig:04_negative_split}. Only 11.5 percent of athletes run the
second four laps faster than the first four, which tells you how rare
genuine restraint is. Those who manage it finish a median eight percentile
places higher in their events than everyone else. Holding back when you
feel strongest goes against every instinct a start line provokes, and it
demonstrably works.

\chart{04_negative_split}{Only 11.5 percent of athletes run the second half
faster, and they finish a median eight percentile places higher.}

And then there is the roxzone, the transition area, where
Figure~\ref{fig:04_roxzone} documents the cheapest time in the sport being
left on the floor. Front-runners spend about 4:22 of their whole race in
transit between stations. The back of the field spends nearly thirteen
minutes. That eight-and-a-half-minute difference involves no exercise at
all: it is the accumulated cost of standing with hands on knees, walking
slowly, and hesitating before picking up the next implement. A mid-pack
athlete who simply moved through the roxzone with purpose would gain more
time than a month of training buys. Almost nobody does it.

\chart{04_roxzone}{Roxzone time by season and by finishing position. The gap
between front and back is about eight minutes of pure transition.}

% ================================================================== 8
\section{What actually makes a fast athlete}

With sixteen timed segments per athlete and a million races, we can ask the
training question directly: what predicts a Hyrox result? The correlation
matrix in Figure~\ref{fig:05_correlation_heatmap} gives the first, slightly
unsettling answer: everything predicts everything. There are no cold cells.
Athletes fast on the row are fast on the sled and fast on the runs. Hyrox
has no niches to hide in; it rewards general fitness and punishes any
weakness, because no strength elsewhere fully pays a weakness's bill.

\chart[width=0.95\linewidth,height=0.58\textheight]{05_correlation_heatmap}{Correlations between all segments. There are
no low-correlation refuges; the sport punishes weakness anywhere.}

Rank the segments by how strongly they predict the finish,
Figure~\ref{fig:05_finish_correlations}, and one bar stands clear of the
rest: total running time correlates 0.91 with the final result, ahead of
every station. Running wins twice over, since it is also the biggest share
of the clock, and the reason is physiological: your running reflects the
aerobic engine that also powers recovery inside and between every station.
If you can train only one quality for this sport, the data's answer is
unambiguous. Train running endurance. Figure~\ref{fig:05_variance_by_segment}
adds the complementary view in raw minutes: running also spreads the field
by the most total time, and among the stations, wall balls and lunges offer
the largest pools of winnable time. Running first, the big grinds second:
the priority order falls straight out of the data.

\chart{05_finish_correlations}{Correlation of each segment with the finish.
Running, at 0.91, is the best predictor by a clear margin.}

\chart{05_variance_by_segment}{The same idea in minutes: the spread of each
segment across the field. Running offers the most time, wall balls and
lunges the most among stations.}

The most interesting chart in this section is the archetype map,
Figure~\ref{fig:05_archetypes}. Every athlete is placed by running ability
on one axis and station strength on the other, with colour showing their
average finishing position. The colour flows along the diagonal, which is
expected: fast and strong beats slow and weak. The revealing part is the two
off-diagonal corners. The pure runner with weak stations and the strong
athlete who cannot run finish at nearly the same middling level, and both
finish behind a balanced athlete of the same total ability.
Figure~\ref{fig:05_imbalance_penalty} reduces this to one line: the more
lopsided an athlete's profile, the worse their average result, steadily and
without exception. Hyrox taxes specialists. Your weakest quality is exposed
eight times per race, and the marginal minute is almost always cheaper to
find in what you are bad at.

\chart{05_archetypes}{Every athlete mapped by running versus station
ability, coloured by average result. Balance beats specialisation at every
level.}

\chart{05_imbalance_penalty}{The specialist tax: average result worsens
steadily as an athlete's run-versus-strength profile gets more lopsided.}

Two more angles complete the picture.
Figure~\ref{fig:05_run_station_ratio} shows that faster athletes spend a
\emph{larger} share of their race running, not because they run more, but
because their station work is so quick that it barely dents the clock; for
the back of the field the stations swallow the race and the runs become
recovery jogs. And Figure~\ref{fig:05_consistency_score} measures how even
each athlete's sixteen segments are relative to the field: the best
athletes are the most uniform, with no segment far out of line. The elite
do not have weapons. They have no weaknesses, which is a different and far
more instructive thing.

\chart{05_run_station_ratio}{Faster athletes spend a larger share of their
race running, because their stations cost so little.}

\chart{05_consistency_score}{Evenness across all sixteen segments against
results. The best athletes have no weak link.}

% ================================================================== 9
\section{The sharp end}

While the middle of the field was being reshaped by growth, the front was
simply getting faster. Figure~\ref{fig:06_winning_time_trend} tracks the
median event-winning time per season: for men it fell from 1:04:44 in
Season 1 to 57:47 in Season 8, with the women's curve falling in parallel.
Remember that this happened while the median athlete drifted slower, and
the apparent contradiction dissolves: the two lines measure different
populations. The median tracks the swelling novice base; the winning time
tracks a professionalising elite, now training specifically for this sport,
racing for prize money and sponsorship.

\chart{06_winning_time_trend}{Median event-winning time per season. Winners
got about seven minutes faster while the median athlete got slower.}

How recently the elite arrived is startling.
Figure~\ref{fig:06_top200_by_season} shows when the 200 fastest times ever
were run: almost all of them in the last two seasons. The record book is
being rewritten in real time, and Tables~\ref{tab:top10_men} and
\ref{tab:top10_women} name the current authors. Alexander Roncevic's 50:38
in Cologne leads the men, Lena Putters' 56:53 in Maastricht the women, and
nearly every entry on both lists comes from Season 7 or 8, most from the
same handful of championship venues where fast fields gather. On the
current trajectory, few of these marks will survive long.
Table~\ref{tab:station_bests} adds the fastest plausible time ever recorded
on each station, using the 99.9th percentile rather than raw minima because
the raw records are contaminated by timing-chip glitches (nobody has done
the burpee broad jumps in nineteen seconds).

\chart{06_top200_by_season}{When the 200 fastest times in history were set:
overwhelmingly in the last two seasons.}

\tabfrag{top10_men}{The ten fastest men's individual finishes on record, one
entry per athlete.}

\tabfrag{top10_women}{The ten fastest women's individual finishes on record,
one entry per athlete.}

\tabfrag{station_bests}{Best plausible segment times in individual racing,
taken at the 99.9th percentile to filter chip-timing glitches.}

Two consequences of a faster front and a broader base follow directly.
Figure~\ref{fig:06_elite_field_gap} shows the distance between the top one
percent and the median athlete widening season after season, the sport
stratifying the way running did as it matured. And
Figure~\ref{fig:06_benchmark_share} resolves an apparent paradox: even as
the absolute number of fast athletes grows, sub-60 men and sub-70 women
shrink as a \emph{share} of the field, because the denominator of newcomers
grows faster still. Percentages and counts can point in opposite directions
during rapid growth, and here they do.

\chart{06_elite_field_gap}{Minutes between the top one percent and the
median athlete, widening every season.}

\chart{06_benchmark_share}{Benchmark finishes as a share of the field. The
elite deepens in absolute terms while diluting as a percentage.}

I want to close this section with what I find the most instructive chart in
the report, Figure~\ref{fig:06_elite_time_budget}. It compares how the ten
fastest men in history distribute their race time against the median male
field, each as a share of their own total. The two bars are almost
identical: about half running, the same station proportions, the same
transition share. The world's best have no secret allocation, no exotic
pacing model, no re-balanced race. They are simply faster everywhere, in
the same shape as everyone else. There is no shortcut to reverse-engineer,
and in a sport this young, that is strangely liberating: the path forward
for everyone is the same, raise the whole profile.

\chart{06_elite_time_budget}{The top ten men all-time and the median field
distribute their race time almost identically. Elites are uniformly faster,
not differently balanced.}

% ================================================================== 10
\section{Careers: what happens after the first race}

Because athletes are tracked by name, the numbers in this section are
careful estimates rather than exact accounting, but the patterns are far too
large to be artifacts. Start with the uncomfortable one:
Figure~\ref{fig:07_races_per_athlete} shows that about 69 percent of
singles athletes appear exactly once. Each further start is roughly half as
common as the previous one. Hyrox is spectacular at recruiting and much
weaker at retaining, and the sport's growth currently runs on a treadmill
of first-timers.

\chart{07_races_per_athlete}{Races per athlete. About seven in ten singles
athletes have raced exactly once.}

The irony is that coming back pays quickly and reliably.
Figure~\ref{fig:07_improvement_curve} and Table~\ref{tab:improvement} track
repeat racers against their own debut: a median 1.3 minutes faster by race
two, three minutes by race four, and 4.5 minutes by race eight, with most
of the gain banked in the first three or four starts. Early improvement is
cheap because a first Hyrox is as much a lesson as a fitness test; simply
knowing how to pace lap one and move through the roxzone is worth minutes.
Figure~\ref{fig:07_career_traj_pct} shows the deeper version: measured in
field position rather than minutes, repeaters keep climbing even after
their raw times level off, and remember they are climbing through a field
whose front end is itself improving. Racecraft compounds.

\chart{07_improvement_curve}{Median improvement versus each athlete's first
race. About 4.5 minutes by race eight, most of it banked early.}

\tabfrag{improvement}{Career improvement in minutes against each athlete's
debut. Negative numbers mean faster.}

\chart{07_career_traj_pct}{The same careers measured in field percentile.
Position keeps improving even after raw times plateau.}

Figure~\ref{fig:07_pb_rate} translates this into the sport's favourite
currency, the personal best, or PB: an athlete's fastest Hyrox finish so
far. More than half of all second races produce one. By race eight, fewer
than a third do. Early PBs are nearly automatic because the bar is a single
previous attempt; late PBs must be earned against your own accumulated
best, which makes a veteran's PB at race eight a rarer and more impressive
thing than the near-guaranteed one at race two.

\chart{07_pb_rate}{The share of athletes setting a personal best (PB, their
fastest Hyrox so far) at each race of their career.}

Retention, Figure~\ref{fig:07_retention} and Table~\ref{tab:retention}, has
settled just under 40 percent season over season, which is genuinely strong
for an event this punishing. But because each incoming cohort is so
enormous, even 40 percent retention leaves every season numerically
dominated by debutants. Both facts are true at once: the sport keeps its
people reasonably well, and it still grows almost entirely by recruiting
new ones. Converting more one-timers into second-timers, with that
near-guaranteed PB as the carrot, is the highest-leverage growth move the
sport has. Table~\ref{tab:most_active} shows the other extreme, the
archive's most prolific racers with dozens of starts, statistical outliers
who form the connective tissue of every local Hyrox community.

\chart{07_retention}{Share of each season's athletes who raced again the
following season, settling just under 40 percent.}

\tabfrag{retention}{Retention by season transition. The Season 8 to 9
transition is excluded because Season 9 has barely begun.}

\tabfrag{most_active}{The most prolific singles racers in the archive.}

Which nations produce the fastest athletes? Figure~\ref{fig:07_nationality_performance}
ranks countries by their median athlete, and the podium is Portugal
(1:21:21), Spain and Ireland rather than the giant markets. The explanation
is composition once more: Germany, Britain and America bring hundreds of
thousands of casual first-timers, while smaller travelling contingents are
self-selected enthusiasts. This chart measures who signs up, not national
fitness. And Figure~\ref{fig:07_nationality_diversity} counts nationalities
per season, from 46 in the first winter to more than 220 now, which is more
than most Olympic sports can claim. Whatever else is true, the reach is
real.

\chart{07_nationality_performance}{Nations ranked by their median athlete's
finish. Small, self-selected contingents beat giant casual markets.}

\chart{07_nationality_diversity}{Athlete nationalities per season: 46 in
Season 1, more than 220 today.}

% ================================================================== 11
\section{Three more questions the splits can answer}

A few analyses did not fit the main narrative but are too good to leave
out. First, ageing by movement. Figure~\ref{fig:08_age_station_heatmap}
shows how much slower each station gets across the age brackets, and the
gradient is a small physiology lesson: the explosive stations, burpee broad
jumps and wall balls, degrade fastest with age, while the erg machines
barely age at all. Power leaves before endurance does. For a masters
athlete the chart doubles as a training plan: defend the explosive stations
that time is quietly taxing, and lean on the aerobic ones where the years
cost almost nothing.

\chart[width=0.95\linewidth,height=0.58\textheight]{08_age_station_heatmap}{How each station slows with age. Explosive
stations age fastest; the ergs barely age at all.}

Second, a shortcut fitness test. Figure~\ref{fig:08_best_run_lap} relates
each athlete's single fastest kilometre of the day to their overall finish,
and the correlation is 0.79. One flat-out kilometre, buried anywhere in the
race, predicts the whole result remarkably well, because it exposes the
same aerobic engine everything else runs on. If you want to estimate your
Hyrox potential without racing one, a hard standalone kilometre is a
legitimate, nearly free proxy.

\chart{08_best_run_lap}{An athlete's best single kilometre predicts their
whole race with a correlation of 0.79.}

Third, how the race changes shape as you move back through the field.
Figure~\ref{fig:08_time_allocation_by_cohort} splits each finishing group's
race into running, stations and transitions: from front to back, the
station and roxzone share grows steadily at running's expense. And
Figure~\ref{fig:08_station_position_effect} runs the cleanest control in
the report, normalising every athlete's stations against their own average,
which removes fitness from the equation entirely. Even then, the late-race
grind stations stand out as relatively slowest for everyone. The second
half of a Hyrox is shaped by accumulated fatigue, not by which movements
anyone happens to be good at, which is one more reason the pacing chapter
matters more than any station-specific trick.

\chart{08_time_allocation_by_cohort}{How each finishing group's race splits
between running, stations and transitions. Slower athletes' races are
swallowed by the stations.}

\chart{08_station_position_effect}{Each athlete's stations normalised
against their own average. Late-race stations are relatively slowest for
everyone; fatigue, not skill, shapes the second half.}

% ================================================================== 12
\section{The rise of Mixed Doubles}

One corner of the doubles boom deserves its own look, with an important
caveat attached from the outset: the scraped archive does not capture every
mixed-gender team, so the numbers here are a partial, indicative view rather
than a complete count. With that firmly in mind,
Figure~\ref{fig:08_mixed_doubles_overview} breaks the doubles field we can
see into its three categories, teams of two men, two women, or one of each.
Among the entries we have, mixed teams make up roughly a quarter of doubles,
with same-sex pairs taking the larger shares. Because the true figure is
almost certainly higher than what the archive records, and because doubles is
now the sport's biggest format overall, mixed racing is clearly a meaningful
slice of Hyrox rather than a niche. The right panel maps where these teams
race, again only among the records we hold: the United States, Germany and
the United Kingdom host the most, with Spain and France close behind, which
points to a broadly Western spread across the big markets. The appeal is easy
to understand whatever the exact totals turn out to be. A mixed team is one
of the most accessible ways into the sport, a couple or a pair of friends
splitting the workload over the same course, and it plausibly helps explain
why doubles has pulled ahead of solo racing.

\chart{08_mixed_doubles_overview}{The doubles field split by team gender
(left) and host countries for mixed teams (right). Mixed-doubles coverage in
the archive is incomplete, so read both panels as indicative rather than
exact.}

% ================================================================== 13
\section{A view from India}

I want to end close to home. India has hosted Hyrox in three cities so far,
Mumbai from Season 7, then Bengaluru and Delhi joining in Season 8, and
these races show both how new and how different the Indian market still is.
Figure~\ref{fig:08_india_vs_world_participation} shows pure growth: Mumbai's debut drew
around 1{,}000 results, and one season later the three Indian venues together
drew about 8{,}000. Add the nationality trend from
Figure~\ref{fig:01_emerging_markets}, Indian athletes going from 100
results in Season 6 to nearly 6{,}000 in Season 8, and India is on the
steepest growth curve of any major market in the sport.

\chart{08_india_vs_world_finish}{Indian venues against the world: finish distributions for Mumbai and Bengaluru versus the global field.}

\chart{08_india_vs_world_participation}{Participation at Indian venues by season.}

Figure~\ref{fig:08_india_vs_world_finish} shows the flip side of that youth. Mumbai's median Open
finish is 1:54:08 and Bengaluru's 2:00:35, against a world median of
1:29:55, which also makes Bengaluru the slowest large venue on earth by
median (Figure~\ref{fig:02_city_difficulty}). It would be easy to misread
that as a verdict. It is a timestamp. A field made almost entirely of
first-timers finishes where first-timers finish everywhere, and every
mature Hyrox market once looked like this. The career section showed what
happens next: second races are 1.3 minutes faster on median and over half
of them are personal bests. As Indian athletes rack up starts, and as the
sport reaches a country full of exactly the young 25 to 34 demographic
that dominates Hyrox globally, those medians will fall fast.

Finally, Figure~\ref{fig:08_india_vs_world_mixed_doubles} hints at a
different flavour to the Indian doubles field: among the entries the archive
records, mixed teams sit a little higher in India than in the world as a
whole, alongside a noticeably larger share of men's pairs. As with the global
mixed-doubles numbers above, coverage of these teams is incomplete, so this is
a suggestive comparison rather than a firm one. Either way, the interesting
question about India is not why its medians are slow today; it is how quickly
this young market will mature.

\chart{08_india_vs_world_mixed_doubles}{Mixed Doubles popularity in India versus the rest of the world.}

% ================================================================== 14
\section{What I take away, and what to be careful about}

Three ideas kept resurfacing no matter which direction I sliced a million
results.

Hyrox is a pacing exam wearing a strongman costume. Half the race is
running, running predicts the result better than anything else, the
biggest gaps between finishing groups appear in fade, steadiness and
transitions, and the world's best allocate their race exactly like everyone
else while executing every part faster. Nothing about winning this sport is
glamorous. The athletes who beat you slowed down less than you did.

The sport is becoming two sports in one, healthily. The front of the field
professionalised, taking seven minutes off winning times in eight seasons,
while the base broadened so fast that medians drifted the other way.
Doubles is on course to overtake solo racing, women are approaching parity,
and the emerging markets, India, China and South Africa above all, are
compounding at rates the mature markets never saw. Europe and the US
dominate today; the map itself says that will not remain the whole story.

And the next frontier is the second race. Seven in ten athletes never come
back, yet returners improve by minutes and PB more than half the time.
No lever in the sport is worth more than converting the first-timer who
just discovered they could finish.

\subsection*{What to be careful about}

Everything above rests on data that was publicly scrapable from Hyrox's
results servers, and that has consequences worth restating plainly. This is
not an official export and Hyrox has not verified any of it. Athletes who
started but did not finish are largely absent, so nothing here should be
read as a completion or dropout rate. Careers are tracked by athlete name
and sex, so same-named athletes occasionally merge and renamed athletes
split, making all career figures estimates. City comparisons mix course,
climate and, above all, field composition, and cannot separate them.
Roxzone times are partly derived where the archive omits them. Mixed-gender
doubles are only partially captured, so every mixed-doubles figure in
Section~12 and the India comparison is a lower bound and indicative view
rather than a complete count. Season 9 was a few weekends old at collection
time and is marked as partial wherever it appears. And the station records
use the 99.9th percentile because the raw minima contain physically
impossible chip glitches. None of these caveats threatens the report's main
conclusions, which rest on patterns spanning hundreds of thousands of
athletes, but precise small numbers deserve a grain of salt.

\appendix

% ================================================================== A
\section{Appendix: reference tables}

The table below is the complete segment reference behind the report: median
times for every station plus running and roxzone totals, by division and
sex, computed over all clean races. Pick your division and sex and every
median becomes a concrete benchmark for your next race.

\tabfrag{station_stats}{Median segment times by division and sex, clean
races, all time.}

\end{document}
